%-----------------------------------------------------------------------
%  Packages, column type and the \hd heading macro.
%  \input by main.tex -- not compiled on its own.
%  Shares no file with ../../nn_surrogate/, so this document can be its own
%  Overleaf project and go to one reader alone.
%-----------------------------------------------------------------------
\usepackage[T1]{fontenc}
\usepackage[utf8]{inputenc}
\usepackage{lmodern}
\usepackage[margin=1in]{geometry}
\usepackage{amsmath,amssymb}
\usepackage{booktabs}
\usepackage{tabularx}
\usepackage{microtype}

\newcolumntype{L}{>{\raggedright\arraybackslash}X}
% Y{w}: an X column with relative width w. Within one table the weights must
% sum to the number of Y/X columns.
\newcolumntype{Y}[1]{>{\hsize=#1\hsize\raggedright\arraybackslash}X}
\newcommand{\code}[1]{\texttt{\small #1}}

% Matches ../../nn_surrogate/master_plan/preamble.tex, so a paragraph moved
% between the two directions does not have to be retyped -- with one deliberate
% improvement: \dg is \ensuremath here, where the master plan writes $^{\circ}$.
% That form is text-mode only, so \dg inside math ($\sim 0.71\dg$) closes the
% surrounding math and reopens it, and the error it throws points at the end of
% the table rather than at the cell. \ensuremath works in both modes.
\newcommand{\kv}{\kappa_{v}}
\newcommand{\dg}{\ensuremath{^{\circ}}}

% \par first: \hd is a run-in heading and must open its own paragraph. Without
% it a heading at the top of an \input file joins the previous file's last
% paragraph, which is invisible in the source and obvious in the PDF.
\newcommand{\hd}[1]{\par\smallskip\noindent\textbf{#1}\ \ignorespaces}

\setlength{\parindent}{0pt}
\setlength{\parskip}{4pt}
\renewcommand{\arraystretch}{1.06}

% Last resort before an overfull line: let TeX stretch rather than overflow.
\emergencystretch=2em
